\chapter{METODE PENELITIAN}
\label{ch:metode}

Bab ini memaparkan metodologi penelitian yang digunakan untuk membangun sistem pelaporan masalah lingkungan berbasis fusi multimodal citra dan teks. Pemaparan diawali dengan diagram alir kerangka berpikir yang menyatukan seluruh tahap penelitian, dilanjutkan dengan analisis kebutuhan, akuisisi dan persiapan dataset Jakarta CRM, lingkungan komputasi pelatihan, pra-proses dan ekstraksi \textit{embedding}, fusi multimodal, algoritma klasifikasi CatBoost dengan \textit{Posterior Gaussian Sampling} (PGS), rancangan eksperimen komparatif, metrik evaluasi, ekspor ONNX, perancangan server FastAPI, perancangan aplikasi Flutter SmartCityApps beserta UML, mekanisme perutean instansi terdekat, dan perancangan layar. Pendekatan fusi awal yang diadopsi merujuk pada metodologi yang dipublikasikan oleh \citet{hanifPenugasanOtomatis2024}.

% =====================================================================
\section{Diagram Alir Kerangka Berpikir}
\label{sec:kerangka}
% =====================================================================

Penelitian ini disusun ke dalam tiga tahap besar yang saling berkesinambungan, yaitu (1) tahap akuisisi dan persiapan data, (2) tahap pelatihan model, dan (3) tahap \textit{deployment}. Diagram alir sederhana pada Gambar~\ref{fig:kerangka-berpikir} merangkum alur keseluruhan dari perolehan data mentah, pelatihan model 9 instansi, \textit{deployment} inferensi pada backend FastAPI, hingga aplikasi Android SmartCityApps yang berinteraksi dengan pengguna akhir sebagai klien.

Seluruh diagram alir dan diagram aktivitas pada bagian ini (Gambar~\ref{fig:kerangka-berpikir}--\ref{fig:kerangka-deployment}) disusun menggunakan notasi \textit{activity diagram} UML \citep{boochUML2005} yang dilengkapi konvensi simbol \textit{flowchart} standar \citep{iso5807}, yaitu \textit{node} awal dan akhir, kotak proses, belah ketupat keputusan, batang \textit{fork/join} untuk percabangan paralel, serta \textit{partition} atau \textit{swimlane} untuk memisahkan tanggung jawab antar aktor atau lingkungan eksekusi. \textit{Partition} inilah yang tampil sebagai kotak pengelompok di dalam diagram (kotak di dalam kotak), dan secara notasi berbeda dari kontainer berlapis pada diagram komponen dan \textit{deployment} di Bagian~\ref{sec:server-design}.

\begin{figure}[htbp]
    \centering
    \includegraphics[width=0.98\textwidth]{bab_3/kerangka_berpikir.png}
    \caption{Diagram alir kerangka berpikir penelitian secara keseluruhan.}
    \label{fig:kerangka-berpikir}
\end{figure}

\subsection{Kerangka Berpikir Tahap Akuisisi Data}

Tahap akuisisi data mencakup proses \textit{web scraping} portal aduan Jakarta CRM, pembersihan dan normalisasi label dinas, pemetaan label mentah ke sembilan kategori target, serta pembagian data secara stratifikasi. Diagram alir tahap ini ditunjukkan pada Gambar~\ref{fig:kerangka-akuisisi}.

\begin{figure}[htbp]
    \centering
    \includegraphics[width=0.85\textwidth, height=0.85\textheight, keepaspectratio]{bab_3/kerangka_akuisisi.png}
    \caption{Diagram alir tahap akuisisi data dari \textit{web scraping} portal CRM Jakarta sampai pembagian data stratifikasi.}
    \label{fig:kerangka-akuisisi}
\end{figure}

\subsection{Kerangka Berpikir Tahap Pelatihan Model}

Tahap pelatihan model dilakukan di lingkungan komputasi awan (\textit{cloud}) RunPod yang diakses melalui SSH. Tahap ini mencakup ekstraksi \textit{embedding} citra dengan empat \textit{encoder} kandidat, ekstraksi \textit{embedding} teks dengan empat \textit{encoder} kandidat, fusi awal melalui konkatenasi vektor ter-normalisasi L2, pelatihan algoritma CatBoost dengan PGS, dan evaluasi pada metrik \textit{macro F1-score}. Karena alur pelatihan terdiri atas dua fase yang panjang, diagram aktivitas dibagi menjadi dua bagian, yaitu fase persiapan (Gambar~\ref{fig:kerangka-pelatihan-setup}) dan fase eksperimen (Gambar~\ref{fig:kerangka-pelatihan-eksperimen}).

\begin{figure}[htbp]
    \centering
    \includegraphics[width=0.85\textwidth, height=0.85\textheight, keepaspectratio]{bab_3/activity_runpod_setup.png}
    \caption{Fase persiapan pelatihan: setup pod RunPod, instalasi dependensi, dan ekstraksi \textit{embedding} citra dan teks ke berkas \textit{cache}.}
    \label{fig:kerangka-pelatihan-setup}
\end{figure}

\begin{figure}[htbp]
    \centering
    \includegraphics[width=0.72\textwidth, height=0.85\textheight, keepaspectratio]{bab_3/activity_runpod_train_export.png}
    \caption{Fase eksperimen: pelatihan \textit{classifier head} CatBoost+PGS untuk 16 pasangan \textit{encoder}, pemilihan model terbaik, dan ekspor ONNX dengan validasi paritas numerik.}
    \label{fig:kerangka-pelatihan-eksperimen}
\end{figure}

\subsection{Kerangka Berpikir Tahap Deployment}

Tahap \textit{deployment} mencakup ekspor model ke format ONNX, pembungkusan ke kontainer Docker, peluncuran server FastAPI, integrasi ke aplikasi Flutter SmartCityApps, dan perutean instansi terdekat berdasarkan koordinat laporan. Diagram alir tahap ini ditunjukkan pada Gambar~\ref{fig:kerangka-deployment}.

\begin{figure}[htbp]
    \centering
    \includegraphics[width=0.55\textwidth, height=0.85\textheight, keepaspectratio]{bab_3/kerangka_deployment.png}
    \caption{Diagram aktivitas tahap \textit{deployment} mencakup ekspor ONNX, server FastAPI untuk inferensi server-side, aplikasi Android SmartCityApps sebagai klien, perutean instansi terdekat, dan pengujian sistem.}
    \label{fig:kerangka-deployment}
\end{figure}

% =====================================================================
\section{Analisis Kebutuhan}
\label{sec:kebutuhan}
% =====================================================================

\subsection{Analisis User}

Pengguna utama sistem ini adalah masyarakat Jakarta yang berinteraksi langsung dengan aplikasi mobile untuk melaporkan masalah lingkungan, beserta operator instansi pemerintah daerah yang menerima laporan setelah perutean otomatis. Karakteristik kedua kelompok pengguna dirangkum sebagai berikut.

\begin{enumerate}
    \item \textbf{Warga Pelapor}: pengguna umum yang menemukan kejadian di lingkungan kota, contohnya jalan rusak, sampah menumpuk, atau pohon tumbang. Mereka memiliki perangkat \textit{smartphone} dengan kamera dan akses GPS, kemampuan menulis narasi singkat dalam Bahasa Indonesia, serta kebutuhan akan respons cepat dari pihak yang berwenang.
    \item \textbf{Operator Instansi}: petugas dinas pemerintah daerah yang menerima laporan setelah dirutekan oleh sistem. Mereka membutuhkan laporan yang sudah terkategorisasi dengan benar, lengkap dengan koordinat lokasi dan bukti visual, agar dapat segera ditindaklanjuti.
\end{enumerate}

\subsubsection{Latar Belakang Institusi}

Penelitian ini berorientasi pada konteks tata kelola Pemerintah Provinsi DKI Jakarta yang memiliki kanal aduan publik berbasis CRM. Struktur instansi yang menjadi target perutean meliputi Dinas Bina Marga, Dinas Perhubungan, Satpol PP, Dinas Pertamanan dan Hutan Kota, Badan Pembina BUMD, Dinas Cipta Karya Tata Ruang dan Pertanahan, Dinas Sumber Daya Air, dan Kelurahan, beserta kategori \textit{Instansi Lain} sebagai kelas \textit{catch-all} untuk laporan di luar kelas utama. Kesembilan instansi target ini ditetapkan secara konsisten pada Bagian~\ref{sec:dataset}.

\subsection{Analisis Aplikasi Sejenis}

Beberapa aplikasi pelaporan publik yang sudah berjalan menjadi pembanding pada penelitian ini, sebagaimana dirangkum pada Tabel~\ref{tab:aplikasi-sejenis}. Aplikasi-aplikasi tersebut secara umum sudah mendukung pelaporan berbasis foto dan teks, namun belum memanfaatkan fusi multimodal antara citra dan teks untuk mengotomatisasi kategorisasi laporan ke instansi yang tepat.

\begin{table}[htbp]
\centering
\caption{Perbandingan aplikasi pelaporan publik sejenis.}
\label{tab:aplikasi-sejenis}
\begin{tabular}{p{2.4cm} p{3.0cm} p{3.0cm} p{3.6cm}}
\toprule
\textbf{Aplikasi} & \textbf{Modalitas Masukan} & \textbf{Klasifikasi Otomatis} & \textbf{Catatan} \\
\midrule
JAKI (DKI Jakarta) & Foto, teks, lokasi & Manual oleh petugas & Kategori dipilih warga \\
LAPOR! (Nasional) & Teks, lampiran & Manual oleh \textit{back office} & Kanal lintas instansi \\
Qlue & Foto, teks, lokasi & Sebagian otomatis & Klasifikasi berbasis aturan \\
Penelitian ini & Foto, teks, lokasi & Otomatis berbasis fusi multimodal & DINOv3 + mE5 + CatBoost+PGS \\
\bottomrule
\end{tabular}
\end{table}

\subsection{Rumusan dan Solusi Kebutuhan}

Setiap pertanyaan penelitian (RQ) yang diajukan pada Bab~\ref{ch:pendahuluan} dipetakan ke komponen solusi yang dirancang dalam metodologi ini. Pemetaan tersebut dirangkum pada Tabel~\ref{tab:rumusan-solusi}.

\begin{table}[htbp]
\centering
\caption{Pemetaan rumusan masalah ke komponen solusi.}
\label{tab:rumusan-solusi}
\setlength{\tabcolsep}{4pt}
\begin{tabular}{p{0.7cm} p{6.3cm} p{6.3cm}}
\toprule
\textbf{RQ} & \textbf{Rumusan Masalah} & \textbf{Solusi yang Dirancang} \\
\midrule
RQ1 & Bagaimana merancang aplikasi pelaporan multimodal? & Aplikasi Flutter dengan layar buat laporan, tinjauan AI, riwayat, dan peta. \\
RQ2 & Bagaimana arsitektur fusi awal yang efektif? & DINOv3 + mE5 + konkatenasi 2048-d + CatBoost+PGS. \\
RQ3 & Bagaimana ekspor ke ONNX dan server inferensi? & \textit{Pipeline} ONNX \textit{opset} 17 + FastAPI multipart + Docker. \\
RQ4 & Bagaimana mekanisme perutean otomatis? & Routing instansi terdekat berdasarkan kelas prediksi 9 instansi dan koordinat latitude-longitude laporan. \\
RQ5 & Mengapa inferensi \textit{on-server} dipilih ketimbang \textit{on-device}? & Analisis ukuran \textit{artifact} (\textit{image encoder} 1,15 GB), latensi, ketersediaan, dan biaya; inferensi dipusatkan pada server FastAPI. \\
\bottomrule
\end{tabular}
\end{table}

% =====================================================================
\section{Akuisisi dan Persiapan Dataset}
\label{sec:dataset}
% =====================================================================

\subsection{Sumber Data CRM Jakarta}

Dataset penelitian diperoleh dari portal aduan publik Jakarta yang berfungsi sebagai sistem CRM (\textit{Citizen Relationship Management}). Setiap laporan disusun dari tiga komponen utama, yaitu (1) gambar lampiran sebagai bukti visual kejadian, (2) narasi tekstual dalam Bahasa Indonesia yang ditulis bebas oleh warga, dan (3) label dinas yang menunjukkan instansi pemerintah daerah penanggung jawab. Volume data mentah hasil pengumpulan adalah \textbf{88.302 laporan} yang terdistribusi pada \textbf{487 nama instansi (SKPD)} berbeda.

\subsection{Struktur Data Multimodal}

Setiap baris dataset memuat sepasang masukan multimodal (gambar dan teks) yang berasal dari satu kejadian yang sama, beserta label dinas sebagai target klasifikasi. Satu contoh baris representatif ditampilkan pada Gambar~\ref{fig:sample-row-1}; dua contoh tambahan dengan karakteristik berbeda (saluran air mampet dengan narasi panjang dan pembatas jalan pudar dengan permohonan pengecatan ulang) disajikan pada Lampiran~\ref{app:sample-multimodal}.

\begin{figure}[htbp]
    \centering
    \includegraphics[width=0.95\textwidth]{sample_row_1.png}
    \caption{Contoh struktur data multimodal: foto ojek parkir di area halte bus, narasi singkat warga, dan label \textit{Dinas Perhubungan}. Tiga kolom yang ditampilkan adalah kolom \textit{Gambar}, kolom \textit{Laporan}, dan kolom \textit{Label}.}
    \label{fig:sample-row-1}
\end{figure}

Tiga kolom inti yang menjadi masukan dan keluaran sistem dirangkum sebagai berikut.

\begin{itemize}
    \item \textbf{Kolom Gambar}: foto kejadian yang diambil oleh warga menggunakan kamera ponsel. Karakteristik visual sangat beragam, mulai dari foto siang hari beresolusi tinggi, foto malam hari dengan pencahayaan rendah, hingga foto sudut tidak biasa dengan oklusi. Resolusi piksel berkisar dari 100 hingga 4.000 piksel pada sisi terlebar.
    \item \textbf{Kolom Laporan}: narasi tekstual yang ditulis bebas oleh warga dalam Bahasa Indonesia. Narasi sering kali pendek dan informal, kadang memuat singkatan, sapaan, atau campuran ragam bahasa daerah. Contohnya, narasi \textit{``ojol parkir di area halte bus tolong dibantu tindaklanjuti laporan ini''} secara semantik merujuk pada pelanggaran parkir kendaraan di area halte bus, sehingga relevan bagi Dinas Perhubungan, sekalipun narasi tidak menyebut nama dinas tersebut.
    \item \textbf{Kolom Label}: nama dinas penanggung jawab yang berasal dari proses pemetaan oleh operator portal CRM. Label ini menjadi target klasifikasi multikelas pada sembilan kategori akhir setelah pemetaan (Bagian~\ref{sec:dataset}).
\end{itemize}

Sifat multimodal dan komplementer kedua modalitas terlihat pada contoh Gambar~\ref{fig:sample-row-1}: foto memperlihatkan kendaraan terparkir di sekitar halte bus, sementara narasi memberi konteks bahwa kendaraan tersebut adalah \textit{ojol} dan tindakannya berstatus pelanggaran. Dua contoh tambahan pada Lampiran~\ref{app:sample-multimodal} memperkuat pola yang sama, yaitu narasi panjang yang menjelaskan dampak banjir dari saluran air mampet (label Kelurahan) serta foto pembatas jalan pudar yang disertai permohonan pengecatan ulang (label Dinas Perhubungan). Ketiga contoh tersebut mengilustrasikan bahwa baik citra maupun teks membawa informasi yang tidak sepenuhnya saling tergantikan, sehingga fusi multimodal diperlukan untuk klasifikasi yang akurat.

\subsection{Metadata Tambahan}

Selain tiga kolom inti di atas, setiap laporan juga menyimpan metadata tambahan yang berguna untuk audit dan analisis, yaitu identitas unik laporan, koordinat lokasi (lintang dan bujur), waktu pengiriman, dan tautan ke berkas gambar. Metadata ini tidak menjadi masukan langsung bagi model klasifikasi, namun digunakan untuk pembagian data, analisis kesalahan, dan rekonstruksi konteks pada tahap evaluasi.

\subsection{Web Scraping}

Pengumpulan data dilakukan dengan teknik \textit{web scraping} berbasis Python sesuai prinsip yang diuraikan oleh \citet{mitchellWebScraping2018}. Skrip \textit{scraping} dikembangkan menggunakan kombinasi pustaka \texttt{requests} untuk pengiriman permintaan HTTP, \texttt{BeautifulSoup} dan \texttt{lxml} untuk penguraian HTML, serta \texttt{Selenium} untuk halaman yang dirender oleh JavaScript secara dinamis. Tahap-tahap utama \textit{scraping} adalah:

\begin{enumerate}
    \item \textbf{Identifikasi sumber}: pemetaan URL halaman daftar laporan dan halaman rincian laporan, beserta pola \textit{pagination}.
    \item \textbf{Pengunduhan}: pengiriman permintaan GET dengan pembatasan laju (\textit{rate limit}) maksimum satu permintaan per detik dan \textit{user-agent} yang mengidentifikasi tujuan akademik.
    \item \textbf{Penguraian}: ekstraksi medan teks laporan, label dinas, tanggal, koordinat, dan tautan gambar.
    \item \textbf{Pengunduhan gambar}: penyimpanan gambar lampiran laporan ke struktur direktori berdasar label dinas.
    \item \textbf{Penyimpanan}: penulisan metadata terstruktur ke berkas CSV dan unduhan gambar ke struktur direktori berbasis label dinas.
    \item \textbf{Penanganan kegagalan}: \textit{retry} dengan \textit{exponential backoff} pada kesalahan jaringan dan \textit{logging} permintaan yang gagal.
\end{enumerate}

\textit{Scraping} menghormati berkas \texttt{robots.txt} dan \textit{Terms of Service} portal sumber, serta tidak mengumpulkan data yang dapat mengidentifikasi individu di luar narasi laporan publik.

\subsection{Pembersihan dan Normalisasi}

Tahap pembersihan dilakukan dengan langkah-langkah berikut:

\begin{itemize}
    \item \textbf{Deduplikasi teks}: penghapusan baris dengan kolom \texttt{laporan} yang persis sama. Tahap ini menghilangkan sekitar 30\% laporan yang merupakan duplikat hasil \textit{copy-paste} (terutama dari kanal Kelurahan).
    \item \textbf{Normalisasi nama dinas}: penyeragaman ejaan dan kapitalisasi nama instansi yang ditulis berbeda-beda pada data sumber.
    \item \textbf{Pembersihan teks}: pemotongan URL via \textit{regex} \texttt{https?://\textbackslash S+|www\textbackslash .\textbackslash S+}, normalisasi unicode ke bentuk NFC, dan \textit{whitespace collapsing}.
    \item \textbf{Validasi gambar}: pengecekan keterbacaan berkas gambar dan penghapusan baris dengan gambar rusak.
\end{itemize}

Setelah pembersihan, dataset bersih berukuran \textbf{61.773 laporan} siap dipakai untuk tahap berikutnya.

\subsection{Pemetaan ke Sembilan Instansi Target}

\paragraph{Definisi terminologi.} Untuk menghindari kerancuan istilah pada bab-bab selanjutnya, ditetapkan tiga definisi berikut yang dipakai secara konsisten di seluruh manuskrip:

\begin{itemize}
    \item \textbf{Instansi (dinas)}: nama lembaga pemerintah daerah yang menjadi target klasifikasi. Sembilan instansi target pada penelitian ini adalah Dinas Bina Marga, Satuan Polisi Pamong Praja (Satpol PP), Dinas Perhubungan, Kelurahan, Dinas Pertamanan dan Hutan, Dinas Sumber Daya Air, Dinas Cipta Karya, Tata Ruang, dan Pertanahan, Badan Pembinaan BUMD, dan Instansi Lain.
    \item \textbf{Instansi prediksi} (\texttt{predicted\_dinas}): keluaran langsung model klasifikasi multimodal berupa salah satu dari sembilan instansi di atas, beserta probabilitas masing-masing kelas.
    \item \textbf{Instansi tujuan terdekat}: hasil tahap routing setelah inferensi. Backend menggunakan \texttt{predicted\_dinas} serta koordinat latitude dan longitude laporan untuk memilih kantor atau unit instansi aktif yang paling dekat dan relevan. Tahap ini \textbf{bukan} agregasi label model; model tetap menghasilkan keluaran pada granularitas sembilan instansi penuh (lihat \texttt{TARGET\_CLASSES} pada \texttt{src/crm/\_\_init\_\_.py}).
\end{itemize}

Istilah ``kategori'' dan ``kelas'' dipakai hanya pada konteks teori klasifikasi (\textit{multi-class problem}, \textit{class weight}, \textit{per-class precision}), bukan sebagai sinonim ``instansi''. Pada bab Hasil dan Pembahasan istilah \textbf{instansi prediksi} digunakan untuk merujuk keluaran model dan \textbf{instansi tujuan terdekat} untuk merujuk hasil perutean berbasis lokasi.

Label dinas pada dataset ini bukan anotasi yang dibuat ulang oleh peneliti, melainkan \textbf{label operasional} (\textit{operational ground truth}) yang diambil dari sistem produksi CRM Jakarta, yaitu instansi yang secara historis benar-benar menindaklanjuti setiap laporan. Penugasan tersebut mengikuti aturan kewenangan resmi pada Surat Keputusan Sekretaris Daerah Provinsi DKI Jakarta Nomor 99 Tahun 2022 yang menetapkan instansi penanggung jawab untuk setiap jenis aduan, sebagaimana didokumentasikan oleh penelitian acuan \citep{hanifPenugasanOtomatis2024}. Dengan demikian, validitas label bertumpu pada keputusan administratif yang dapat ditelusuri ke dasar hukum, bukan pada penilaian subjektif anotator, sehingga re-anotasi manual terhadap 61.773 sampel tidak diperlukan dan masalah \textit{inter-annotator agreement} yang lazim pada pelabelan manual tidak berlaku. Sebagai konsekuensinya, kualitas label mewarisi derajat \textit{noise} operasional yang melekat pada sistem, yaitu potensi ketidakkonsistenan keputusan petugas antar-waktu serta keberadaan laporan yang secara sah dapat ditangani lebih dari satu instansi; karakteristik ini diperlakukan secara eksplisit sebagai ketidakpastian \textit{aleatoric} dan ditinjau pada analisis kesalahan di Bab~\ref{ch:hasil}.

Karena 487 SKPD mentah terlalu \textit{fine-grained} untuk klasifikasi \textit{multi-class} dengan jumlah sampel yang tidak seimbang, label dinas mentah dipetakan ke sembilan instansi target yang lebih representatif dan operasional. Pemetaan ini bersifat deterministik dan mengikuti dua prinsip: (1) \textit{penggabungan hierarkis}, yaitu suku dinas dan unit pelaksana teknis dikonsolidasikan ke dinas induknya sesuai struktur organisasi Pemerintah Provinsi DKI Jakarta; dan (2) \textit{pemotongan berbasis frekuensi}, yaitu delapan instansi dengan volume laporan terbanyak (mencakup lebih dari 90 persen total laporan) dipertahankan sebagai kelas tersendiri, sedangkan ekor-panjang yang heterogen dikonsolidasikan ke satu kelas \textit{catch-all} ``Instansi Lain''. Strategi dua-prinsip ini selaras dengan penelitian acuan \citep{hanifPenugasanOtomatis2024}. Pemetaan dan distribusi sampel per instansi dirangkum pada Tabel~\ref{tab:kategori-target}.

\begin{table}[htbp]
\centering
\caption{Distribusi sampel per instansi target setelah pemetaan.}
\label{tab:kategori-target}
\begin{tabular}{p{0.7cm} p{6.5cm} r r}
\toprule
\textbf{ID} & \textbf{Kategori (Dinas)} & \textbf{Jumlah} & \textbf{Persen} \\
\midrule
0 & Dinas Bina Marga & 13.525 & 21,9\% \\
1 & Dinas Perhubungan (termasuk UP Perparkiran, UPSL) & 11.678 & 18,9\% \\
2 & Kelurahan & 11.980 & 19,4\% \\
3 & Satuan Polisi Pamong Praja & 8.405 & 13,6\% \\
4 & Dinas Pertamanan dan Hutan Kota & 6.000 & 9,7\% \\
5 & Dinas Sumber Daya Air & 2.605 & 4,2\% \\
6 & Badan Pembina BUMD (TransJakarta, Pasar Jaya, PAM Jaya) & 1.684 & 2,7\% \\
7 & Dinas Cipta Karya Tata Ruang dan Pertanahan & 1.475 & 2,4\% \\
8 & Instansi Lain (\textit{catch-all}) & 4.421 & 7,2\% \\
\midrule
& \textbf{Total} & \textbf{61.773} & \textbf{100\%} \\
\bottomrule
\end{tabular}
\end{table}

\subsection{Pembagian Data Stratifikasi}

Dataset bersih dibagi ke dalam tiga partisi (\textit{train}, \textit{validation}, \textit{test}) dengan rasio \textbf{70:15:15}, menggunakan fungsi \texttt{StratifiedShuffleSplit} dari pustaka \texttt{scikit-learn} untuk menjaga proporsi setiap kelas pada masing-masing partisi. Hasil pembagian dirangkum pada Tabel~\ref{tab:split}, dan diverifikasi bahwa selisih proporsi kelas antar partisi tidak melampaui 0{,}01\%.

\begin{table}[htbp]
\centering
\caption{Pembagian dataset stratifikasi.}
\label{tab:split}
\begin{tabular}{l r r}
\toprule
\textbf{Partisi} & \textbf{Jumlah Sampel} & \textbf{Persentase} \\
\midrule
Train & 43.241 & 70\% \\
Validation & 9.266 & 15\% \\
Test & 9.266 & 15\% \\
\midrule
\textbf{Total} & \textbf{61.773} & \textbf{100\%} \\
\bottomrule
\end{tabular}
\end{table}

Setiap baris pada partisi diberi anotasi \texttt{row\_id} unik agar \textit{cache embedding} dapat disejajarkan kembali dengan label saat pelatihan algoritma klasifikasi.

\subsection{Penanganan Ketidakseimbangan Kelas: Pembobotan, Bukan \textit{Resampling}}
\label{sec:bobot-kelas}

Karena distribusi kelas tidak seimbang, pelatihan algoritma CatBoost menggunakan \textit{class weights} berdasarkan formula yang juga dipakai oleh \citet{hanifPenugasanOtomatis2024}:

\begin{equation}
    w_l = \frac{N}{C \cdot N_l},
    \label{eq:classweight}
\end{equation}
dengan $N$ total sampel pelatihan, $C = 9$ jumlah kelas, dan $N_l$ jumlah sampel kelas ke-$l$. Bobot ini dipakai sebagai parameter bobot kelas pada CatBoost untuk mengompensasi kelas minoritas.

Penelitian ini secara sengaja \textbf{tidak} menerapkan \textit{resampling} (baik \textit{oversampling} seperti SMOTE maupun \textit{undersampling}) untuk menyeimbangkan distribusi, melainkan memilih kompensasi berbasis biaya melalui bobot kelas pada Persamaan~\ref{eq:classweight}. Keputusan ini berlandaskan tiga pertimbangan yang didukung literatur mutakhir.

Pertama, manfaat \textit{balancing} menurun drastis pada \textit{classifier} kuat. \citet{elorToSMOTE2022} mengevaluasi SMOTE beserta variannya pada 73 dataset dan menyimpulkan bahwa \textit{balancing} tidak meningkatkan kinerja prediksi untuk \textit{classifier} kuat seperti \textit{gradient-boosted trees}; manfaatnya hanya nyata pada \textit{learner} lemah dan cenderung lenyap setelah pemilihan metrik dan penyetelan \textit{hyperparameter} yang tepat. Karena \textit{classifier head} pada penelitian ini adalah CatBoost (\textit{boosted trees}) yang menyediakan bobot kelas secara \textit{native}, pembobotan biaya merupakan jalur yang lebih hemat dan tepat-sasaran dibandingkan sintesis sampel.

Kedua, SMOTE bekerja buruk pada ruang fitur berdimensi tinggi. Masukan \textit{classifier} pada penelitian ini adalah vektor fusi berdimensi 2048 hasil konkatenasi \textit{embedding} DINOv3 dan multilingual-E5. \citet{maldonadoFWSMOTE2022} menunjukkan bahwa interpolasi SMOTE yang bergantung pada jarak Euclidean menjadi tidak bermakna pada dimensi tinggi karena jarak antar-instance cenderung konvergen ke nilai yang sama (\textit{curse of dimensionality}), sehingga sampel sintetis yang dihasilkan dapat menyimpang dari \textit{manifold} kelas minoritas dan justru menambah \textit{noise}. Menyintesis sampel pada ruang \textit{embedding} padat semacam ini berisiko menurunkan kualitas keputusan alih-alih memperbaikinya.

Ketiga, konsistensi metodologis dengan baseline dan domain. Penelitian acuan \citet{hanifPenugasanOtomatis2024} yang diperluas pada studi ini juga memilih pembobotan kelas, dan karya sebidang pada laporan warga berbahasa Indonesia \citep{zahroFairerPublicComplaint2025} memilih pendekatan berbasis fungsi rugi (\textit{focal loss}) ketimbang \textit{resampling}. Dengan mengadopsi pembobotan kelas, hasil penelitian ini tetap sebanding secara metodologis sekaligus menghindari risiko kebocoran informasi (\textit{leakage}) yang dapat timbul bila sampel sintetis dari kelas minoritas bocor lintas partisi \textit{train}/\textit{validation}. Sebagai konsekuensi, distribusi sampel pada Tabel~\ref{tab:split} tetap timpang apa adanya; ketimpangan tersebut dikompensasi pada tingkat fungsi tujuan, bukan dihapus dari data, dan dampaknya dimitigasi lebih lanjut melalui pemilihan \textit{macro F1} sebagai metrik utama (Bagian~\ref{sec:skenario-pengujian}).

% =====================================================================
\section{Lingkungan Komputasi: Pelatihan di Cloud via RunPod}
\label{sec:runpod}
% =====================================================================

\subsection{Justifikasi Cloud Compute}

Ekstraksi \textit{embedding} dari empat \textit{encoder} citra besar dan empat \textit{encoder} teks besar pada 61.773 sampel membutuhkan komputasi GPU yang berkapasitas. Sebagai contoh, \textit{forward pass} EVA-02 \textit{Large} pada masukan 448x448 piksel memerlukan memori GPU yang melampaui 12 GB pada \textit{batch size} efektif. Eksekusi pada perangkat lokal kelas \textit{laptop} tidak memungkinkan, sehingga seluruh pelatihan dilakukan pada layanan komputasi awan RunPod yang menyediakan akses langsung ke GPU NVIDIA modern.

\subsection{Konfigurasi Pod RunPod}

Konfigurasi \textit{pod} yang digunakan dirangkum pada Tabel~\ref{tab:runpod}. Pemilihan RTX 4090 didasarkan pada kombinasi VRAM 24 GB yang memadai untuk batch ekstraksi \textit{embedding} dan biaya per jam yang ekonomis dibanding kelas A100.

\begin{table}[htbp]
\centering
\caption{Konfigurasi \textit{pod} RunPod untuk pelatihan dan ekstraksi \textit{embedding}.}
\label{tab:runpod}
\begin{tabular}{p{4.5cm} p{8cm}}
\toprule
\textbf{Komponen} & \textbf{Spesifikasi} \\
\midrule
GPU & NVIDIA GeForce RTX 4090 (24 GB VRAM) \\
CPU & vCPU 8 \textit{core} (AMD EPYC) \\
RAM & 32 GB \\
\textit{Container image} & \texttt{runpod/\allowbreak{}pytorch:\allowbreak{}2.1.0-\allowbreak{}py3.10-\allowbreak{}cuda12.1.1} \\
Penyimpanan & 50 GB \textit{persistent volume} \\
Mode jaringan & SSH \textit{public key authentication} \\
Port \textit{exposed} & 22 (SSH), 8888 (Jupyter) \\
\bottomrule
\end{tabular}
\end{table}

\subsection{Workflow Pelatihan dan Manajemen Lingkungan}

Akses ke \textit{pod} dilakukan melalui kunci publik SSH yang didaftarkan pada akun RunPod. \textit{Workflow} pelatihan mencakup sinkronisasi dataset dengan \texttt{rsync}, pemasangan dependensi (\texttt{torch}, \texttt{transformers}, \texttt{timm}, \texttt{catboost}, \texttt{onnxruntime-gpu}), eksekusi notebook melalui \textit{SSH tunnel} ke Jupyter Lab, penyimpanan \textit{artifact} pada \textit{persistent volume}, dan pengunduhan hasil ke mesin lokal. Lingkungan direproduksi dari berkas konfigurasi dengan versi dependensi presisi, dan \textit{forward pass} \textit{encoder} memakai presisi campuran FP16 dengan ukuran \textit{batch} yang disesuaikan per \textit{encoder} (misalnya 32 untuk DINOv3 pada 224x224 dan 8 untuk EVA-02 pada 448x448). Konteks visual alur kerja ini sudah disajikan pada Gambar~\ref{fig:kerangka-pelatihan-setup}, sedangkan rincian langkah SSH serta praktik manajemen lingkungan secara penuh disajikan pada Lampiran~\ref{app:runpod-workflow}.

% =====================================================================
\section{Pemrosesan, Ekstraksi \textit{Embedding}, dan Fusi Dini}
\label{sec:preprocessing}
% =====================================================================

Tahap pemrosesan data multimodal, ekstraksi \textit{embedding}, dan fusi dini menyiapkan masukan tunggal berdimensi 2048 untuk \textit{classifier head} CatBoost. Pada sisi citra, setiap gambar dibaca dan diorientasikan ulang sesuai EXIF, kemudian di-\textit{resize} dengan \textit{square padding} ke 224x224 (atau 448x448 untuk EVA-02) mengikuti \textit{AutoImageProcessor} bawaan DINOv3, lalu dinormalisasi per kanal dengan rerata serta simpangan baku ImageNet sesuai Persamaan~\ref{eq:norm}. Augmentasi tidak diterapkan karena \textit{encoder} dibekukan. Pada sisi teks, narasi dipangkas spasinya dan dideduplikasi (penghapusan teks identik pada tahap pembersihan dataset), dipotong ke 256 token (mencakup sebagian besar laporan tanpa pemotongan), diberi prefiks instruksi untuk varian mE5-Instruct, ditokenisasi dengan tokenizer bawaan tiap \textit{encoder}, lalu \textit{embedding} keluarannya dinormalisasi L2.

Ekstraksi \textit{embedding} dilakukan dengan empat kandidat \textit{encoder} citra (DINOv2-L, DINOv3-ViT-L, Hiera-L, EVA-02-L) dan empat kandidat \textit{encoder} teks (mE5-Large-Instruct \texttt{intfloat/multilingual-e5-large-instruct}, BGE-M3, IndoBERT-Large-p2, Cendol-mT5-Large), seluruhnya menghasilkan vektor berdimensi 1024. \textit{Encoder} teks yang dipakai pada konfigurasi terbaik adalah varian \textbf{\textit{instruct}} dari \textit{multilingual-E5 large} (selanjutnya disingkat mE5), yang diberi prefiks instruksi pada setiap masukan. Inferensi dilakukan secara \textit{batch} dengan presisi campuran FP16 dan hasilnya disimpan pada \textit{cache} terkompresi sehingga eksperimen ablasi pasangan \textit{encoder} dapat dilakukan tanpa menjalankan ulang \textit{encoder}.

Skema fusi yang digunakan adalah konkatenasi langsung dua vektor \textit{embedding} yang sudah dinormalisasi L2 per modalitas, mengikuti pendekatan \citet{hanifPenugasanOtomatis2024}, menghasilkan vektor fusi berdimensi 2048 yang menjadi masukan tunggal CatBoost. Reduksi dimensi (PCA) tidak diterapkan karena CatBoost berbasis pohon \textit{oblivious} mampu menangani fitur padat berdimensi tinggi serta untuk menghindari hilangnya informasi diskriminatif. Rincian lengkap berupa urutan langkah pemrosesan, tabel kandidat \textit{encoder}, formula normalisasi L2, dan persamaan vektor fusi disajikan pada Lampiran~\ref{app:preprocessing-embedding}.

% =====================================================================
\section{Algoritma Klasifikasi: CatBoost dengan PGS}
\label{sec:catboost}
% =====================================================================

\subsection{Justifikasi Pemilihan CatBoost sebagai \textit{Classifier Head}}
\label{sec:justifikasi-catboost}

Pemilihan CatBoost sebagai \textit{classifier head} pada \textit{pipeline} fusi awal penelitian ini didasarkan pada delapan properti yang selaras dengan kebutuhan teknis dan operasional sistem, yaitu: (1) \textit{ordered boosting} \citep{prokhorenkovaCatBoost2018} yang memitigasi \textit{target leakage} sehingga generalisasi pada dataset berukuran sedang (61.773 sampel) lebih stabil; (2) pohon simetris (\textit{oblivious trees}) yang membuat inferensi cepat dan ramah \textit{cache} CPU pada server FastAPI; (3) ketahanan empiris pada masukan tabular berdimensi tinggi seperti vektor fusi 2048 dimensi hasil konkatenasi \textit{embedding} DINOv3 dan \textit{multilingual-E5}; (4) kalibrasi probabilitas melalui Posterior Gaussian Sampling \citep{ustimenkoCatBoostUncertainty2023} yang sekaligus memberikan ukuran ketidakpastian untuk antarmuka pengguna; (5) pelatihan deterministik dengan \textit{random seed} eksplisit sehingga ablasi dapat direproduksi dan diaudit; (6) jalur ekspor resmi ke ONNX dengan paritas numerik tinggi terhadap implementasi \textit{native}, selaras dengan \textit{deployment} server FastAPI; (7) beban komputasi pelatihan yang ringan (satu \textit{classifier head} selesai dalam hitungan menit), sehingga enam belas kombinasi ablasi dapat dijalankan dalam satu sesi karena kedua \textit{encoder} dibekukan; dan (8) penanganan fitur kategoris \textit{native} via \textit{ordered target statistics} untuk perluasan masa depan, misalnya menambahkan fitur lokasi atau waktu pelaporan tanpa mengganti algoritma. Penelitian ini menggunakan CatBoost sebagai satu-satunya algoritma klasifikasi; eksplorasi konfigurasi (CatBoost dengan dan tanpa PGS) didokumentasikan pada Bagian~\ref{sec:eksperimen}.

\subsection{Hyperparameter}

Hyperparameter CatBoost yang dipakai dirangkum pada Tabel~\ref{tab:catboost-hp}.

\begin{table}[htbp]
\centering
\caption{Hyperparameter CatBoost.}
\label{tab:catboost-hp}
\begin{tabular}{p{5.0cm} p{7.0cm}}
\toprule
\textbf{Parameter} & \textbf{Nilai} \\
\midrule
Jumlah iterasi & 1500 \\
Kedalaman pohon & 6 \\
\textit{Learning rate} & 0,05 \\
Fungsi tujuan & \textit{MultiClass logloss} \\
Regularisasi L2 daun & 3,0 \\
Bobot kelas & sesuai Persamaan~\ref{eq:classweight} \\
\textit{Posterior sampling} & Aktif \\
\textit{Early stopping rounds} & 50 (pada validasi) \\
\textit{Task type} & CPU (atau GPU bila tersedia) \\
\textit{Random seed} & 42 \\
\bottomrule
\end{tabular}
\end{table}

\subsection{\textit{Posterior Gaussian Sampling} (PGS)}

PGS adalah teknik kuantifikasi ketidakpastian pada \textit{gradient boosting} yang diperkenalkan oleh \citet{ustimenkoCatBoostUncertainty2023}. Saat pelatihan dengan opsi \textit{posterior sampling} aktif, CatBoost membangun \textit{virtual ensembles} yang dapat dipanggil saat inferensi dengan jumlah anggota ensembel virtual sebanyak 30. Setiap anggota menghasilkan distribusi probabilitas \textit{softmax} per kelas. Probabilitas akhir adalah rerata posterior:

\begin{equation}
    \bar{p}_k = \frac{1}{M} \sum_{m=1}^{M} p_k^{(m)},
    \label{eq:pgs-mean}
\end{equation}
dengan $M = 30$ anggota ensembel dan $p_k^{(m)}$ probabilitas kelas $k$ dari anggota ensembel ke-$m$.

Ketidakpastian \textit{epistemic} diestimasi dengan \textit{mutual information} antara prediksi dan parameter model, sebagaimana selisih \textit{entropy of mean} dan \textit{mean of entropy}:

\begin{equation}
    \mathcal{U}_{\text{epistemic}} = H(\bar{\mathbf{p}}) - \frac{1}{M} \sum_{m=1}^{M} H(\mathbf{p}^{(m)}),
    \label{eq:pgs-mi}
\end{equation}
dengan $H(\mathbf{p}) = -\sum_k p_k \log p_k$ adalah entropi Shannon. Nilai $\mathcal{U}_{\text{epistemic}}$ yang tinggi menandakan model tidak yakin terhadap prediksi karena ketidakpastian parameter, sehingga laporan dengan nilai tinggi dapat dirutekan ke tinjauan manual.

\subsection{Pelatihan dengan Validation Set dan Early Stopping}

Pelatihan dilakukan pada partisi \textit{train} (43.241 sampel) dengan validasi pada partisi \textit{validation} (9.266 sampel). Mekanisme \textit{early stopping} dengan ambang 50 iterasi memerintahkan pelatihan berhenti apabila \textit{macro F1-score} pada set validasi tidak meningkat selama 50 iterasi berturut-turut, sehingga risiko \textit{overfitting} tertekan dan jumlah pohon optimal ditemukan secara otomatis.

\subsection{Pemilihan Model Terbaik}

Model terbaik dipilih berdasarkan \textit{macro F1-score} pada partisi validasi. \textit{Checkpoint} terbaik disimpan dalam format \textit{native} CatBoost dan kemudian diekspor ke ONNX untuk \textit{deployment}.

% =====================================================================
\section{Eksperimen Komparatif (Ablasi)}
\label{sec:eksperimen}
% =====================================================================

Bagian ini mendokumentasikan rancangan eksperimen komparatif yang dijalankan untuk memilih kombinasi \textit{encoder} terbaik dan untuk mengevaluasi kontribusi PGS pada algoritma CatBoost.

\subsection{Ablasi Matriks Multimodal}

Penelitian ini menjalankan eksperimen matriks penuh, yaitu seluruh pasangan kombinasi antara empat \textit{encoder} citra (DINOv2-L, DINOv3-ViT-L, Hiera-L, EVA-02-L) dan empat \textit{encoder} teks (mE5-Large-Instruct, BGE-M3, IndoBERT-Large-p2, Cendol-mT5-Large), sehingga total 16 kombinasi. Setiap kombinasi dilatih dengan konfigurasi CatBoost yang identik. Tujuan ablasi matriks ini adalah memetakan kontribusi setiap \textit{encoder} terhadap akurasi akhir, sekaligus memilih pasangan terbaik untuk \textit{deployment}.

\subsection{Ablasi PGS On/Off}

Pengaruh \textit{Posterior Gaussian Sampling} dievaluasi dengan menjalankan setiap kombinasi pada matriks ablasi dengan dua varian, yaitu CatBoost \textit{baseline} tanpa PGS dan CatBoost dengan PGS aktif beserta 30 anggota ensembel virtual saat inferensi. Untuk setiap pasangan dicatat \textit{accuracy}, \textit{macro F1-score}, selisih relatif ($\Delta$F1, $\Delta$Acc), dan rerata \textit{epistemic uncertainty}. Tujuan ablasi ini adalah memverifikasi bahwa PGS memberikan peningkatan yang konsisten lintas kombinasi \textit{encoder}, di samping memberikan keluaran ketidakpastian yang dapat diaudit.

\subsection{Tabel Rancangan Eksperimen}

Matriks eksperimen lengkap dirangkum pada Tabel~\ref{tab:matrix-eksperimen}. Setiap baris mewakili satu kombinasi (\textit{encoder} citra, \textit{encoder} teks) yang dievaluasi dua kali (\textit{baseline} dan PGS).

\begin{table}[htbp]
\centering
\caption{Matriks rancangan eksperimen ablasi (4 \textit{encoder} citra \(\times\) 4 \textit{encoder} teks \(\times\) 2 varian \textit{head}).}
\label{tab:matrix-eksperimen}
\small
\begin{tabular}{c p{3.2cm} p{4.0cm} p{3.2cm}}
\toprule
\textbf{ID} & \textbf{Encoder Citra} & \textbf{Encoder Teks} & \textbf{Head} \\
\midrule
E01 & DINOv2-L & mE5-Large-Instruct & CatBoost / CatBoost+PGS \\
E02 & DINOv3-ViT-L & mE5-Large-Instruct & CatBoost / CatBoost+PGS \\
E03 & Hiera-L & mE5-Large-Instruct & CatBoost / CatBoost+PGS \\
E04 & EVA-02-L & mE5-Large-Instruct & CatBoost / CatBoost+PGS \\
E05 & DINOv2-L & BGE-M3 & CatBoost / CatBoost+PGS \\
E06 & DINOv3-ViT-L & BGE-M3 & CatBoost / CatBoost+PGS \\
E07 & Hiera-L & BGE-M3 & CatBoost / CatBoost+PGS \\
E08 & EVA-02-L & BGE-M3 & CatBoost / CatBoost+PGS \\
E09 & DINOv2-L & IndoBERT-Large-p2 & CatBoost / CatBoost+PGS \\
E10 & DINOv3-ViT-L & IndoBERT-Large-p2 & CatBoost / CatBoost+PGS \\
E11 & Hiera-L & IndoBERT-Large-p2 & CatBoost / CatBoost+PGS \\
E12 & EVA-02-L & IndoBERT-Large-p2 & CatBoost / CatBoost+PGS \\
E13 & DINOv2-L & Cendol-mT5-Large & CatBoost / CatBoost+PGS \\
E14 & DINOv3-ViT-L & Cendol-mT5-Large & CatBoost / CatBoost+PGS \\
E15 & Hiera-L & Cendol-mT5-Large & CatBoost / CatBoost+PGS \\
E16 & EVA-02-L & Cendol-mT5-Large & CatBoost / CatBoost+PGS \\
\bottomrule
\end{tabular}
\end{table}

% =====================================================================
\section{Metrik Evaluasi}
\label{sec:metrics-method}
% =====================================================================

Metrik utama dan pendukung yang dipakai untuk mengevaluasi setiap eksperimen sudah didefinisikan secara teoretis pada Bagian~\ref{sec:metrics}. Pada metodologi ini metrik yang dilaporkan untuk setiap eksperimen adalah:

\begin{enumerate}
    \item \textbf{Macro F1-score}, sebagaimana Persamaan~\ref{eq:macrof1}, sebagai metrik utama pemilihan model.
    \item \textbf{Akurasi} dan \textbf{Balanced Accuracy} sebagai indikator umum.
    \item \textbf{Per-Class F1, Precision, Recall} untuk setiap dari sembilan kategori, agar bias terhadap kelas mayoritas teridentifikasi.
    \item \textbf{Confusion Matrix} berukuran $9 \times 9$, divisualisasikan sebagai \textit{heatmap}.
    \item \textbf{Macro ROC-AUC} dengan skema \textit{one-vs-rest}.
    \item \textbf{Analisis ketidakpastian}: distribusi $\mathcal{U}_{\text{epistemic}}$ pada partisi uji, dan korelasi antara nilai ketidakpastian dengan kesalahan prediksi.
\end{enumerate}

Keseluruhan alur inferensi yang menghasilkan metrik-metrik di atas, mulai dari citra dan teks hingga prediksi sembilan instansi, \textit{confidence}, dan routing instansi terdekat, ditunjukkan pada Gambar~\ref{fig:flow-inferensi}. \textit{Confusion matrix} aktual berukuran $9 \times 9$ pada konfigurasi terbaik disajikan pada Bab~\ref{ch:hasil}.

\begin{figure}[htbp]
    \centering
    \includegraphics[width=0.98\textwidth, height=0.55\textheight, keepaspectratio]{bab_3/flow_inferensi.png}
    \caption{Alur inferensi multimodal server-side dari citra dan teks ke prediksi 9 instansi, confidence, dan routing instansi terdekat berdasarkan koordinat.}
    \label{fig:flow-inferensi}
\end{figure}

% =====================================================================
\section{Pipeline Ekspor ONNX}
\label{sec:onnx-method}
% =====================================================================

\textit{Pipeline} ekspor menghasilkan tiga \textit{artifact} ONNX terpisah yang dimuat oleh server FastAPI saat \textit{startup}.

\subsection{Ekspor Encoder Citra}

\textit{Encoder} citra terbaik diekspor melalui kombinasi \texttt{torch.jit.trace} dan \texttt{torch.onnx.export} dengan \textit{opset} versi 17. Penggunaan \texttt{torch.jit.trace} diperlukan karena beberapa varian \textit{encoder}, khususnya DINOv3, belum sepenuhnya didukung oleh \textit{exporter} otomatis. Berkas hasil disimpan beserta konfigurasi pra-proses (\textit{mean}, \textit{std}, ukuran masukan) dalam berkas pendamping \textit{JSON}.

\subsection{Ekspor Encoder Teks}

\textit{Encoder} teks terbaik diekspor menggunakan utilitas ekspor ONNX yang menangani modul \textit{attention mask} dan SDPA secara benar. Bobot model yang melampaui ukuran maksimum berkas tunggal ONNX disimpan sebagai bobot eksternal. Tokenizer dikemas sebagai \textit{artifact} pendamping pada direktori yang sama dengan model.

\subsection{Ekspor CatBoost}

Bobot CatBoost diekspor menggunakan metode \textit{native} \texttt{model.save\_model(\allowbreak{}"classifier.onnx",\allowbreak{} format="onnx")} yang menulis \textit{operator domain} \texttt{ai.catboost}. Inferensi pada saat produksi dijalankan dengan \texttt{Execution Provider} CPU karena \textit{operator} CatBoost belum memiliki kernel CUDA pada \textit{ONNX Runtime}.

\subsection{Prosedur Validasi Paritas Numerik}

Validasi paritas numerik dirancang untuk memastikan bahwa model ONNX hasil ekspor menghasilkan keluaran yang konsisten dengan model PyTorch sumber. Prosedurnya adalah sebagai berikut.

\begin{enumerate}
    \item Pilih satu \textit{batch} sampel dari partisi uji sebagai masukan referensi.
    \item Jalankan model PyTorch sumber pada \textit{batch} tersebut dan simpan vektor keluaran sebagai \textit{baseline}.
    \item Jalankan model ONNX pada \textit{batch} yang sama melalui \textit{ONNX Runtime} dengan \textit{Execution Provider} yang sesuai.
    \item Hitung selisih absolut maksimum (\textit{max abs diff}) per komponen vektor keluaran antara kedua model.
    \item Tetapkan ambang batas $10^{-5}$ sebagai kriteria lulus, yang memastikan prediksi top-1 identik antara kedua jalur.
    \item Catat hasil paritas untuk setiap komponen, yaitu \textit{encoder} citra, \textit{encoder} teks, dan algoritma CatBoost.
\end{enumerate}

Laporan numerik hasil validasi paritas akan disajikan pada Bab Hasil dan Pembahasan.

% =====================================================================
\section{Perancangan Sistem Server (FastAPI)}
\label{sec:server-design}
% =====================================================================

\subsection{Arsitektur Server}

Server inferensi membungkus tiga sesi \textit{ONNX Runtime} (\textit{encoder} citra, \textit{encoder} teks, algoritma CatBoost) yang dimuat sekali pada saat \textit{startup}, kemudian menyediakan metode prediksi yang menjalankan \textit{pipeline} fusi awal pada masukan gabungan citra dan teks. Aplikasi Android SmartCityApps berperan sebagai \textit{client} untuk mengirimkan laporan, foto, serta koordinat latitude dan longitude. Proses inferensi model dilakukan pada sisi server melalui backend FastAPI, bukan pada perangkat Android. Diagram komponen dan \textit{deployment} server ditunjukkan pada Gambar~\ref{fig:server-arch}. Diagram tersebut mengikuti notasi \textit{component} dan \textit{deployment diagram} UML \citep{boochUML2005} yang dipadukan dengan pendekatan kontainer berlapis ala \textit{C4 model} \citep{brownC42018}, yaitu sebuah \textit{node} lingkungan eksekusi (Docker Container) yang memuat modul dan komponen di dalamnya (notasi kotak di dalam kotak).

Model yang digunakan dapat disimpan sebagai artefak pada server atau pada \textit{model registry} seperti HuggingFace. Namun, HuggingFace pada rancangan ini diperlakukan sebagai penyimpanan atau registri artefak model, bukan sebagai layanan yang menjalankan inferensi. Dengan pemisahan ini, pembaruan model dapat dilakukan pada sisi server tanpa mewajibkan pengguna memperbarui aplikasi Android.

\begin{figure}[htbp]
    \centering
    \includegraphics[width=0.99\textwidth, height=0.55\textheight, keepaspectratio]{bab_3/deployment_component_drawio.png}
    \caption{Diagram komponen dan deployment SmartCityApps: klien Android, backend FastAPI untuk inferensi server-side, HuggingFace sebagai model registry/storage, dan Supabase.}
    \label{fig:server-arch}
\end{figure}

\subsection{Endpoint \texttt{/health} dan \texttt{/predict}}

Server menyediakan tiga \textit{endpoint}, yaitu \texttt{/health} (GET) sebagai \textit{health check} ringan, \texttt{/predict} (POST \textit{multipart}) sebagai \textit{endpoint} klasifikasi, dan \texttt{/reload-agencies} (POST) untuk memuat ulang registri instansi dari basis data ke \textit{cache} tanpa me-restart server. Endpoint \texttt{/predict} menerima foto laporan, teks laporan, latitude, dan longitude. Spesifikasi \textit{endpoint} dirangkum pada Tabel~\ref{tab:endpoints}.

\begin{table}[htbp]
\centering
\caption{Spesifikasi \textit{endpoint} server FastAPI.}
\label{tab:endpoints}
\begin{tabular}{p{1.7cm} p{2.0cm} p{4cm} p{4.7cm}}
\toprule
\textbf{Path} & \textbf{Method} & \textbf{Masukan} & \textbf{Keluaran} \\
\midrule
\texttt{/health} & GET & (tidak ada) & JSON \texttt{\{status, encoders, gpu, agencies\_source, agencies\_count\}} \\
\texttt{/predict} & POST & \textit{multipart}: \texttt{image} (file), \texttt{laporan} (string), \texttt{latitude}, \texttt{longitude} & JSON \texttt{PredictResponse} berisi prediksi, confidence, probabilitas per kelas, dan assignment instansi terdekat \\
\texttt{/reload-agencies} & POST & (opsional) header \texttt{X-Reload-Token} & JSON \texttt{\{source, count\}}; memuat ulang registri instansi dari Supabase ke \textit{cache} \\
\bottomrule
\end{tabular}
\end{table}

\subsection{Skema Request dan Response}

Skema respons \texttt{PredictResponse} dideklarasikan sebagai kelas Pydantic dengan medan: \texttt{predicted\_dinas} (string), \texttt{predicted\_dinas\_id} (int), \texttt{confidence} (float), \texttt{all\_probabilities} (peta instansi ke skor probabilitas), \texttt{model\_name}, \texttt{model\_version}, \texttt{confidence\_threshold}, \texttt{review\_required}, dan \texttt{assignment}. Medan \texttt{assignment} memuat \texttt{agency\_id}, \texttt{agency\_name}, \texttt{agency\_category}, \texttt{distance\_meters}, dan \texttt{routing\_method}. Validasi otomatis dijamin oleh FastAPI sehingga klien Flutter dapat mengandalkan struktur respons.

\subsection{Output Inferensi dan Confidence Level}

Output proses inferensi terdiri atas kelas prediksi, nilai \textit{confidence}, dan skor probabilitas untuk setiap dari sembilan instansi target. Kelas prediksi menunjukkan instansi yang paling mungkin bertanggung jawab terhadap laporan, sedangkan \textit{confidence} menunjukkan tingkat keyakinan model terhadap kelas tersebut. Nilai ini dihitung dari probabilitas kelas tertinggi pada keluaran model multiclass.

Nilai \textit{confidence} digunakan sebagai indikator pendukung pengambilan keputusan. Apabila nilai \textit{confidence} berada di atas ambang yang ditentukan, laporan dapat diarahkan otomatis ke instansi terdekat yang relevan dengan kelas prediksi. Apabila \textit{confidence} berada di bawah ambang, laporan tetap disimpan bersama skor probabilitasnya, tetapi ditandai untuk validasi manual oleh petugas agar risiko kesalahan klasifikasi tidak disembunyikan.

\subsection{Registri Instansi, Perutean Lokasi, dan Strategi Sumber Data}
\label{subsec:registri-instansi}

Keluaran klasifikasi berupa satu dari sembilan kategori instansi, sedangkan penugasan operasional membutuhkan kantor fisik yang konkret beserta koordinatnya. Kedua hal ini berada pada tingkat granularitas yang berbeda: \emph{kategori} bersifat logis (sembilan kelas keluaran model sesuai \texttt{TARGET\_CLASSES}), sementara \emph{kantor} bersifat fisik dan dapat berjumlah lebih dari satu untuk tiap kategori. Sebagai contoh, kategori \textit{Dinas Bina Marga} dilayani oleh kantor dinas tingkat provinsi sekaligus suku dinas pada beberapa kota administrasi. Oleh karena itu sembilan kategori dipetakan ke 16 entri kantor pada registri instansi sebagai relasi satu-ke-banyak; angka 16 bukan perluasan ruang kelas model, melainkan pencacahan lokasi kantor tujuan yang menjadi kandidat perutean.

Perutean berlangsung dua tahap. Tahap pertama, model memilih kategori \(\hat{y} = \arg\max_{k \in \{1,\dots,9\}} P(y = k \mid x)\). Tahap kedua, di antara seluruh kantor aktif yang berkategori \(\hat{y}\), sistem memilih kantor dengan jarak \textit{haversine} terkecil terhadap koordinat laporan \(\ell_{\text{lap}} = (\textit{lat}, \textit{lon})\):

\begin{equation}
    a^{\star} = \arg\min_{a \in \mathcal{A}_{\hat{y}}} d_{\text{haversine}}\!\left(\ell_{\text{lap}}, \ell_a\right),
    \label{eq:nearest-agency}
\end{equation}

\noindent dengan \(\mathcal{A}_{\hat{y}}\) himpunan kantor aktif berkategori \(\hat{y}\) dan \(\ell_a\) koordinat kantor \(a\). Apabila \(\mathcal{A}_{\hat{y}} = \emptyset\), sistem jatuh ke kategori \textit{Instansi lain} sebagai unit peninjauan manual sehingga laporan tidak pernah gagal dirutekan.

Pengelolaan data kantor memiliki tiga opsi rancangan: (A)~\textit{hardcode} pada kode server, (B)~berkas konfigurasi yang dimuat saat \textit{startup}, dan (C)~tabel basis data sebagai \textit{system of record}. Opsi A melanggar prinsip pemisahan data dari logika---mengubah koordinat menuntut \textit{deploy} ulang kode---dan menyulitkan penelusuran asal-usul data (\textit{provenance}). Opsi B memperbaiki kedua hal tersebut, namun belum memungkinkan pengelolaan oleh administrator non-pengembang saat \textit{runtime}. Penelitian ini mengadopsi Opsi C: tabel \texttt{agencies} pada Supabase menjadi \textit{system of record} ternormalisasi yang dirujuk \texttt{report\_assignments.agency\_id} melalui \textit{foreign key}, dengan hak tulis dibatasi peran super admin melalui kebijakan RLS \texttt{agencies\_write\_super\_admin}.

Agar server inferensi tidak bergantung pada ketersediaan basis data pada setiap permintaan, registri dibaca \emph{satu kali} saat \textit{startup} lalu disimpan pada \textit{cache} memori; permintaan \texttt{/predict} berikutnya melakukan perutean murni dari \textit{cache} tanpa kueri basis data sehingga latensi prediksi tidak terdampak. Apabila basis data tidak terjangkau, server secara otomatis jatuh (\textit{graceful degradation}) ke berkas \texttt{agencies\_seed.json} yang dibundel bersama server sehingga perutean tetap berfungsi. Berkas yang sama berperan ganda: sebagai \textit{seed} awal tabel \texttt{agencies} dan sebagai jembatan pemetaan antara nilai \textit{enum} \texttt{issue\_category} (\textit{slug}) di basis data dengan label kategori yang dipakai algoritma perutean, sehingga kode Python, klien Flutter (enum \texttt{IssueCategory}), dan basis data merujuk pada satu definisi instansi yang konsisten (\textit{single source of truth}).

Perubahan data oleh super admin dilakukan melalui antarmuka pengelolaan registri pada aplikasi Flutter, kemudian dipublikasikan ke mesin perutean melalui \textit{endpoint} \texttt{/reload-agencies} yang memuat ulang \textit{cache} server tanpa \textit{restart}; alur lengkap pengelolaannya ditunjukkan pada Gambar~\ref{fig:activity-agency-registry}. Dengan demikian rancangan ini menggabungkan keunggulan Opsi C (ternormalisasi, dapat dikelola administrator, ber-\textit{provenance}) dengan ketahanan operasional berkat \textit{fallback} berbasis berkas. Koordinat tiap kantor diperoleh melalui \textit{geocoding} alamat resmi (OpenStreetMap Nominatim): empat dinas tingkat provinsi (Bina Marga, Perhubungan, Sumber Daya Air, dan Cipta Karya) berkantor di Gedung Dinas Teknis Jl.\ Taman Jatibaru; Satuan Polisi Pamong Praja dan Badan Pembinaan BUMD di kompleks Balai Kota; sedangkan suku dinas tingkat kota administrasi diproksikan ke kompleks Kantor Walikota kotamadya tempatnya berkedudukan, dengan \textit{provenance} per kantor dicatat pada metadata \texttt{agencies\_seed.json}.

\begin{figure}[htbp]
    \centering
    \includegraphics[width=0.95\textwidth, height=0.5\textheight, keepaspectratio]{bab_3/activity_agency_registry.png}
    \caption{\textit{Activity diagram} pengelolaan registri instansi oleh super admin: penyuntingan data dengan penegakan RLS \texttt{agencies\_write\_super\_admin}, lalu pemuatan ulang \textit{cache} server inferensi melalui \texttt{/reload-agencies} dengan \textit{fallback} ke \texttt{agencies\_seed.json} bila basis data tidak terjangkau.}
    \label{fig:activity-agency-registry}
\end{figure}

\subsection{Dockerization}

Server dikemas dalam kontainer Docker yang mengikuti pola \textit{multi-stage build}. \textit{Base image} berbasis CUDA 12.1 dengan \textit{cuDNN}, lapisan instalasi Python 3.11 dan dependensi inferensi (\texttt{onnxruntime-gpu}, \texttt{fastapi}, \texttt{uvicorn}, \texttt{python-multipart}), serta lapisan penyalinan kode dan \textit{artifact} ONNX. \textit{Entrypoint} menjalankan \texttt{uvicorn} pada \textit{host} \texttt{0.0.0.0} \textit{port} 8000 dengan satu \textit{worker}. Berkas Docker Compose mendefinisikan layanan \textit{server} dengan reservasi GPU dan \textit{volume read-only} untuk \textit{artifact} model.

% =====================================================================
\section{Perancangan Skema Basis Data (Supabase)}
\label{sec:db-schema}
% =====================================================================

Basis data \textit{cloud} aplikasi diimplementasikan pada Supabase yang berbasis PostgreSQL. Skema dirancang ternormalisasi dengan tabel \texttt{reports} sebagai akar agregat (\textit{aggregate root}) dan memindahkan konsep bernilai jamak serta historis ke tabel anak khusus, yaitu bukti penyelesaian (\texttt{report\_evidence}), riwayat penugasan instansi (\texttt{report\_assignments}), audit status (\texttt{report\_history}), dan tenggat layanan (\texttt{report\_sla}); registri instansi tujuan disimpan pada tabel \texttt{agencies}, sedangkan profil pengguna pada tabel \texttt{profiles}. Skema ER lengkap ditunjukkan pada Gambar~\ref{fig:db-schema}.

\begin{figure}[htbp]
    \centering
    \includegraphics[width=0.95\textwidth]{bab_3/erd_report_prediction_assignment.png}
    \caption{Skema basis data Supabase ternormalisasi: akar agregat \texttt{reports} dengan tabel anak bukti, penugasan, audit, dan SLA, serta proyeksi-baca yang disinkronkan \textit{trigger} untuk klien \textit{Realtime}.}
    \label{fig:db-schema}
\end{figure}

\subsection{Ringkasan Tabel dan Kontrol Akses}

Skema basis data terdiri atas tujuh tabel dan enam \textit{enum} pendukung. Tabel \texttt{profiles} terhubung satu-ke-satu dengan \texttt{auth.users} milik Supabase Auth dan menyimpan nama, kontak, peran (\texttt{user\_role}), serta instansi penugasan operator. Tabel \texttt{reports} adalah akar agregat yang menyimpan identitas pelapor, narasi, koordinat, keluaran model (\texttt{ai\_prediction}, \texttt{ai\_confidence}, \texttt{ai\_probabilities} dalam format \texttt{jsonb}, dan metadata model), status siklus hidup (\texttt{report\_status}), visibilitas (\texttt{report\_visibility}), serta penanda duplikasi (\texttt{duplicate\_of}). Konsep bernilai jamak dipindahkan ke tabel anak: \texttt{report\_evidence} menampung banyak berkas bukti per laporan (foto awal maupun foto penyelesaian), \texttt{report\_assignments} merekam riwayat penugasan dan penugasan-ulang instansi, \texttt{report\_history} mencatat audit perubahan status, dan \texttt{report\_sla} menyimpan tenggat layanan; tabel \texttt{agencies} menjadi registri instansi tujuan yang dirujuk oleh penugasan. Kebijakan \textit{Row Level Security} (RLS) memastikan pelapor hanya melihat laporan miliknya, petugas hanya mengakses kategori sesuai \texttt{assigned\_agency}, dan administrator memiliki akses lintas instansi. Rincian kolom seluruh tabel, daftar lengkap nilai \textit{enum}, serta kebijakan RLS disajikan pada Lampiran~\ref{app:db-schema-detail}.

Tabel \texttt{agencies} berperan sebagai registri instansi tujuan perutean (\textit{system of record}) yang dirujuk oleh \texttt{report\_assignments.agency\_id} melalui \textit{foreign key}, dan hak tulisnya dibatasi pada peran super admin melalui kebijakan RLS \texttt{agencies\_write\_super\_admin}. Strategi sumber data registri, mekanisme pembacaannya oleh server inferensi (termasuk \textit{caching} dan \textit{fallback}), pemetaan kategori-ke-kantor, serta \textit{provenance} koordinat dibahas tersendiri pada Sub-bab~\ref{subsec:registri-instansi}.

Pemisahan tabel anak ini menghindari rancangan yang menumpuk seluruh konsep pada satu tabel sehingga fitur produksi seperti bukti ganda, penugasan-ulang, eskalasi, penautan laporan duplikat, dan pemantauan SLA dapat ditambahkan tanpa membongkar skema. Agar klien Flutter yang membaca melalui Supabase \textit{Realtime} (\textit{stream} satu tabel \texttt{reports}) tetap bekerja, kolom datar pada \texttt{reports} dipertahankan sebagai \textit{proyeksi-baca} yang disinkronkan otomatis oleh \textit{trigger} \texttt{SECURITY DEFINER}: setiap tulisan pada \texttt{reports} diproyeksikan ke \texttt{report\_evidence} dan \texttt{report\_assignments}, sementara tabel anak menjadi sumber kebenaran historis. Pola \textit{Command Query Responsibility Segregation} ringan ini menyeimbangkan normalisasi dengan kebutuhan pembacaan \textit{realtime} berbasis tabel.

Berkas bukti dan foto profil disimpan pada \textit{bucket} privat Supabase Storage. Mengikuti praktik produksi, basis data tidak menyimpan URL final melainkan pasangan \textit{bucket} dan \texttt{object\_path} pada \texttt{report\_evidence}; URL akses dibuat sebagai \textit{signed URL} berumur singkat saat berkas hendak ditampilkan. Pendekatan ini mencegah tautan yang dapat kedaluwarsa atau bocor mengekspos foto warga yang berpotensi memuat data pribadi, dan memisahkan referensi penyimpanan secara logis (bukan \textit{foreign key} fisik) dari basis data relasional.

Otentikasi pengguna ditangani oleh Supabase Auth dengan metode surel dan kata sandi. Proses registrasi memicu verifikasi surel, yaitu tautan konfirmasi dikirim ke alamat pengguna dan sesi baru menjadi aktif setelah verifikasi, sedangkan pemulihan kata sandi dilakukan melalui tautan \textit{reset} yang juga dikirim melalui surel. Pengiriman seluruh surel transaksional tersebut difasilitasi oleh layanan SMTP yang terhubung dengan Supabase Auth. Pada prototipe ini digunakan SMTP bawaan Supabase yang memiliki batas laju pengiriman sehingga memadai untuk pengujian, dan dapat diganti dengan penyedia SMTP khusus (misalnya SendGrid atau Resend) pada penerapan produksi tanpa mengubah kode aplikasi.

\subsection{Bukti Penyelesaian Laporan}

Pada tahap penyelesaian laporan, petugas diwajibkan mengunggah foto bukti penyelesaian ketika status laporan diubah menjadi \textit{resolved}/selesai. Foto bukti ini menunjukkan bahwa permasalahan telah ditangani di lapangan. Selain foto, petugas dapat menambahkan catatan penyelesaian yang disimpan pada kolom \texttt{resolution\_note}.

Validasi bukti penyelesaian diterapkan pada alur status. Transisi dari \textit{in\_progress} ke \textit{resolved} tidak dijalankan apabila \texttt{resolution\_evidence\_photo\_url} belum tersedia. Metadata \texttt{resolved\_by} dan \texttt{resolved\_at} juga disimpan agar proses penyelesaian dapat diaudit.

% =====================================================================
\section{Perancangan Sistem Aplikasi Mobile (Flutter)}
\label{sec:flutter-design}
% =====================================================================

Aplikasi Flutter dikembangkan dengan arsitektur berbasis fitur (\textit{feature-based architecture}). Dokumentasi rancangan sistem disusun mengikuti notasi UML standar yang diuraikan pada Bagian~\ref{sec:uml}.

\subsection{Use Case Diagram}

Diagram \textit{use case} memetakan interaksi antara aktor utama (Warga Pelapor dan Petugas Instansi) dengan fungsionalitas SmartCityApps, backend FastAPI, dan Supabase. Diagram \textit{use case} ditunjukkan pada Gambar~\ref{fig:usecase-flutter}.

\begin{figure}[htbp]
    \centering
    \includegraphics[width=0.92\textwidth, height=0.88\textheight, keepaspectratio]{bab_3/use_case_smartcityapps.png}
    \caption{\textit{Use case diagram} SmartCityApps dengan prediksi server-side, confidence level, routing instansi terdekat, dan bukti penyelesaian.}
    \label{fig:usecase-flutter}
\end{figure}

\subsection{Use Case Description}

Setiap \textit{use case} utama dirangkum pada Tabel~\ref{tab:uc-ringkas}. Deskripsi lengkap setiap \textit{use case} (alur utama, alur alternatif, pra-kondisi, dan pasca-kondisi) disajikan pada Lampiran~\ref{app:usecase-detail}.

\begin{table}[htbp]
\centering
\caption{Ringkasan \textit{use case} utama SmartCityApps.}
\label{tab:uc-ringkas}
\begin{tabular}{p{4.0cm} p{2.8cm} p{5.2cm}}
\toprule
\textbf{Use Case} & \textbf{Aktor} & \textbf{Deskripsi Singkat} \\
\midrule
Mendaftar/Masuk Akun & Warga Pelapor & Registrasi atau login untuk mengakses fitur pelaporan. \\
Mengambil Foto Laporan & Warga Pelapor & Ambil foto kamera atau galeri sebagai bukti visual. \\
Menulis Narasi Laporan & Warga Pelapor & Tulis narasi Bahasa Indonesia pelengkap foto. \\
Mengirim Laporan & Warga Pelapor, Server & Kirim foto, narasi, dan koordinat ke FastAPI untuk klasifikasi. \\
Meninjau Hasil AI & Warga Pelapor & Lihat instansi prediksi, confidence, dan instansi terdekat. \\
Mengoreksi Kategori & Warga Pelapor & Ganti instansi prediksi secara manual bila keliru. \\
Menyimpan Laporan & Warga Pelapor & Simpan laporan dan hasil klasifikasi ke Supabase. \\
Melihat Riwayat/Peta & Warga Pelapor & Lihat daftar dan lokasi laporan tersimpan. \\
\bottomrule
\end{tabular}
\end{table}

\subsection{Activity Diagram}

\textit{Activity diagram} memvisualisasikan alur aktivitas pada skenario utama pelaporan oleh Warga Pelapor, mulai dari pengambilan foto, inferensi server-side, penyimpanan ke Supabase, sampai kewajiban bukti foto ketika petugas menandai laporan selesai. Gambar~\ref{fig:activity-buat-laporan} menunjukkan alur lengkap tersebut, termasuk percabangan koreksi instansi secara manual pada layar tinjauan AI.

\begin{figure}[htbp]
    \centering
    \includegraphics[width=0.72\textwidth, height=0.88\textheight, keepaspectratio]{bab_3/activity_report_lifecycle.png}
    \caption{\textit{Activity diagram} siklus laporan SmartCityApps dari pelaporan, inferensi server-side, routing instansi terdekat, sampai bukti penyelesaian.}
    \label{fig:activity-buat-laporan}
\end{figure}

\subsection{Class Diagram}

\textit{Class diagram} memetakan kelas-kelas utama pada aplikasi Flutter dan disusun mengikuti praktik desain produksi, bukan sekadar daftar kelas. Tiga keputusan desain ditonjolkan. Pertama, \textit{interface segregation}: kelas \texttt{ReportsRepository} merealisasikan dua antarmuka terpisah, yaitu \texttt{ReportsReader} (operasi baca berbasis \textit{stream} untuk warga dan admin) dan \texttt{ReportAdminActions} (operasi mutasi untuk moderator dan \textit{super-admin}), sehingga layar yang hanya membaca tidak terekspos pada operasi yang mengubah keadaan. Kedua, \textit{strategy pattern} pada hierarki \texttt{ReportsScope} yang bersifat \textit{sealed} (\texttt{CitizenScope}, \texttt{AgencyScope}, \texttt{GlobalScope}) untuk menentukan cakupan laporan yang terlihat berdasarkan peran pengguna yang juga dimodelkan sebagai hierarki \textit{sealed} \texttt{UserRole}. Ketiga, pemisahan tegas antara \textit{service} klasifikasi awan (\texttt{CloudClassificationService} beserta \textit{data transfer object} \texttt{CloudPrediction}) dan model domain (\texttt{CityReport}, \texttt{AiPrediction}, \texttt{ReportAssignment}) yang dipetakan satu-ke-satu ke tabel Supabase. Abstraksi \texttt{AiClassificationService} memiliki dua implementasi, yaitu \texttt{CloudAiClassificationService} (jalur produksi) dan \texttt{DemoAiClassificationService} (\textit{stub} deterministik untuk mode pengujian). Karena cakupannya luas, \textit{class diagram} dipecah menjadi lima bagian fokus yang saling melengkapi agar setiap kelas tetap terbaca pada ukuran cetak: (1) lapisan layanan inti pada Gambar~\ref{fig:class-diagram}, (2) modul otentikasi dan peran pengguna pada Gambar~\ref{fig:class-diagram-auth}, (3) modul laporan dengan repositori dan strategi cakupan pada Gambar~\ref{fig:class-diagram-reports}, (4) model domain dan enum kategori instansi pada Gambar~\ref{fig:class-diagram-domain}, serta (5) modul administrasi registri instansi pada Gambar~\ref{fig:class-diagram-admin}. Enum \texttt{IssueCategory} (sembilan kelas instansi target) menjadi penghubung antarbagian dan dirinci penuh pada Gambar~\ref{fig:class-diagram-domain}.

\begin{figure}[htbp]
    \centering
    \includegraphics[width=\textwidth, keepaspectratio]{bab_3/class_diagram_1.png}
    \caption{\textit{Class diagram} bagian 1 dari 5, lapisan layanan inti: abstraksi \texttt{AiClassificationService} dengan implementasi \textit{cloud} dan \textit{demo}, klien \texttt{CloudClassificationService} beserta DTO \texttt{CloudPrediction}, serta layanan lokasi dan media.}
    \label{fig:class-diagram}
\end{figure}

\begin{figure}[htbp]
    \centering
    \includegraphics[width=\textwidth, keepaspectratio]{bab_3/class_diagram_2.png}
    \caption{\textit{Class diagram} bagian 2 dari 5, modul otentikasi: \texttt{AuthRepository} berbasis Supabase, model \texttt{UserProfile}, dan hierarki \textit{sealed} \texttt{UserRole} dengan tiga peran (\texttt{Citizen}, \texttt{AgencyAdmin}, \texttt{SuperAdmin}).}
    \label{fig:class-diagram-auth}
\end{figure}

\begin{figure}[htbp]
    \centering
    \includegraphics[width=\textwidth, keepaspectratio]{bab_3/class_diagram_3.png}
    \caption{\textit{Class diagram} bagian 3 dari 5, modul laporan: \texttt{CreateReportController}, \texttt{ReportsRepository} dengan antarmuka tersegregasi \texttt{ReportsReader}/\texttt{ReportAdminActions}, dan strategi cakupan \textit{sealed} \texttt{ReportsScope} per peran pengguna.}
    \label{fig:class-diagram-reports}
\end{figure}

\begin{figure}[htbp]
    \centering
    \includegraphics[width=\textwidth, keepaspectratio]{bab_3/class_diagram_4.png}
    \caption{\textit{Class diagram} bagian 4 dari 5, model domain: \texttt{CityReport}, \texttt{AiPrediction}, \texttt{ReportAssignment}, \texttt{ReportHistoryEntry}, \texttt{ReportLocationData}, serta enum \texttt{IssueCategory} (sembilan kelas) dan \texttt{ReportStatus}.}
    \label{fig:class-diagram-domain}
\end{figure}

\begin{figure}[htbp]
    \centering
    \includegraphics[width=0.9\textwidth, keepaspectratio]{bab_3/class_diagram_5.png}
    \caption{\textit{Class diagram} bagian 5 dari 5, modul administrasi registri instansi: \texttt{AgencyRegistryScreen} dan \texttt{AgencyRegistryRepository} (operasi CRUD ke tabel \texttt{agencies} serta pemuatan ulang \textit{cache} server melalui \texttt{/reload-agencies}), beserta model \texttt{Agency} dan \texttt{AgencyReloadResult}. Penulisan dibatasi peran super admin melalui RLS \texttt{agencies\_write\_super\_admin}.}
    \label{fig:class-diagram-admin}
\end{figure}

\subsection{Sequence Diagram}

\textit{Sequence diagram} pada Gambar~\ref{fig:sequence-diagram} menggambarkan urutan pesan antar komponen pada alur klasifikasi \textit{end-to-end}, mulai dari layar buat laporan, backend FastAPI, routing instansi terdekat, penyimpanan Supabase, hingga unggah bukti penyelesaian oleh petugas.

\begin{figure}[htbp]
    \centering
    \includegraphics[width=0.92\textwidth, height=0.85\textheight, keepaspectratio]{bab_3/sequence_report_submission.png}
    \caption{\textit{Sequence diagram} pengiriman laporan, inferensi server-side, penyimpanan assignment, dan bukti penyelesaian.}
    \label{fig:sequence-diagram}
\end{figure}

\subsection{Modul dan Layanan Flutter}

Struktur modul aplikasi Flutter dirangkum pada Tabel~\ref{tab:flutter-modules}.

\begin{table}[htbp]
\centering
\caption{Modul dan layanan utama aplikasi Flutter.}
\label{tab:flutter-modules}
\begin{tabular}{p{3.5cm} p{9cm}}
\toprule
\textbf{Modul} & \textbf{Deskripsi} \\
\midrule
Modul Otentikasi & Pendaftaran, masuk, dan persistensi sesi lokal. \\
Modul Splash & Layar \textit{splash} dan inisialisasi awal aplikasi. \\
Modul Onboarding & Panduan singkat untuk pengguna baru. \\
Modul Dashboard & Ringkasan laporan dan navigasi utama. \\
Modul Reports & Buat laporan, tinjauan AI, riwayat, peta, perutean instansi terdekat, dan bukti penyelesaian. \\
Modul Profile & Pengelolaan profil pengguna. \\
Modul Settings & Konfigurasi aplikasi. \\
Layanan Klasifikasi Awan & Klien HTTP \textit{multipart} yang berkomunikasi dengan server FastAPI. \\
Layanan Lokasi & Akuisisi koordinat GPS dan \textit{reverse geocoding}. \\
Layanan Media & Pemilih kamera dan galeri foto. \\
Layanan Repositori Laporan & Operasi CRUD pada tabel \texttt{reports} dan \texttt{report\_history} Supabase, termasuk penyimpanan prediksi, assignment, bukti penyelesaian, dan \textit{stream} realtime per peran pengguna. \\
Layanan Otentikasi Supabase & Pendaftaran, masuk, keluar, dan pembaruan profil pengguna dengan tiga peran. \\
\bottomrule
\end{tabular}
\end{table}

% =====================================================================
\section{Mekanisme Perutean Instansi}
\label{sec:routing}
% =====================================================================

\subsection{Keluaran Model Tetap Sembilan Instansi}

Model klasifikasi tidak menghasilkan empat kluster operasional. Keluaran model tetap berupa prediksi multiclass pada sembilan instansi target, yaitu satu \texttt{predicted\_dinas}, nilai \textit{confidence}, dan skor probabilitas untuk seluruh kelas. Istilah kluster operasional pada rancangan awal hanya berfungsi sebagai panduan tampilan atau pengelompokan naratif, bukan sebagai label model dan bukan sebagai hasil akhir routing pada arsitektur Phase 2.

Dengan demikian, alur yang dipakai pada implementasi adalah: notebook melatih dan mengevaluasi model 9 kelas, artefak terbaik diekspor ke ONNX, backend FastAPI memuat artefak tersebut, lalu SmartCityApps menerima hasil prediksi dari server. Aplikasi Android tidak menjalankan model secara lokal.

\subsection{Routing Instansi Terdekat Berbasis Lokasi}

Setelah backend memperoleh \texttt{predicted\_dinas}, sistem menentukan instansi tujuan berdasarkan dua masukan: kelas prediksi dan koordinat laporan. Koordinat latitude dan longitude laporan dibandingkan dengan daftar instansi aktif yang memiliki kategori sesuai. Instansi dengan jarak terdekat dipilih sebagai tujuan laporan.

Perhitungan jarak menggunakan formula Haversine karena koordinat berada pada permukaan bumi. Untuk dua titik koordinat \((lat_1, lon_1)\) dan \((lat_2, lon_2)\), jarak permukaan bumi dihitung dengan pendekatan:

\begin{equation}
    d = 2R \cdot \arctan2(\sqrt{a}, \sqrt{1-a}),
\end{equation}

dengan \(R\) adalah jari-jari bumi dan \(a\) dihitung dari selisih lintang serta bujur. Hasil routing yang dikembalikan backend meliputi \texttt{agency\_id}, \texttt{agency\_name}, \texttt{agency\_category}, \texttt{distance\_meters}, dan \texttt{routing\_method}.

Perlu ditegaskan bahwa daftar instansi beserta koordinatnya pada prototipe ini merupakan data contoh representatif (\textit{seeded}) yang ditanamkan pada layanan, bukan registri resmi yang memuat seluruh unit instansi se-DKI Jakarta. Selain itu, perutean berbasis jarak Haversine merupakan penyederhanaan; pada praktik pemerintahan, penugasan laporan umumnya mengikuti batas wilayah administratif (kelurahan atau kecamatan) yang menaungi lokasi kejadian, bukan semata kantor yang terdekat secara garis lurus. Pada penerapan produksi, koordinat instansi idealnya diisi melalui registrasi lokasi mandiri oleh masing-masing Operator Instansi atau diambil dari registri resmi pemerintah; mekanisme tersebut dibahas sebagai arah pengembangan lanjutan pada Bab~\ref{ch:simpulan}.

\subsection{Kasus Confidence Rendah dan Fallback}

Apabila nilai \textit{confidence} berada di bawah ambang yang dikonfigurasi, laporan tetap dapat disimpan bersama seluruh skor probabilitasnya, tetapi ditandai sebagai membutuhkan validasi manual. Apabila tidak ditemukan instansi aktif yang cocok dengan kelas prediksi, backend menggunakan fallback ke unit tinjauan manual agar permintaan tidak gagal secara diam-diam.

% =====================================================================
\section{Perancangan Layar Aplikasi}
\label{sec:layar}
% =====================================================================

Bagian ini menjelaskan secara ringkas fungsi utama, alur interaksi pengguna, dan keterkaitan tiap layar dengan kebutuhan sistem yang dijabarkan pada Bagian~\ref{sec:kebutuhan}. Mockup tampilan antarmuka untuk seluruh layar disajikan terpusat pada Lampiran~\ref{app:ui-mockup} agar pemaparan metode tetap padat dan rancangan visual dapat ditinjau secara berdampingan.

\subsection{Layar Splash dan Onboarding}

Layar \textit{splash} menampilkan logo aplikasi dan indikator pemuatan saat sumber daya awal (sesi otentikasi Supabase, konfigurasi GoRouter, dan registrasi \textit{provider} Riverpod) dimuat. Setelah pemuatan selesai, pengguna baru diarahkan ke alur \textit{onboarding} tiga halaman yang menjelaskan tiga proposisi nilai aplikasi, yaitu pelaporan multimodal, perutean otomatis, dan transparansi status laporan. Tampilan layar dapat dilihat pada Lampiran~\ref{app:ui-mockup}.

\subsection{Layar Buat Laporan}

Layar buat laporan adalah titik masuk utama Warga Pelapor. Layar ini menyediakan tombol ambil foto, \textit{preview} foto, medan teks narasi multibaris, indikator status GPS, dan tombol klasifikasi yang aktif ketika gambar, deskripsi, latitude, dan longitude tersedia. Data tersebut dikirim ke backend FastAPI; SmartCityApps tidak menjalankan model secara lokal. Tampilan layar dapat dilihat pada Lampiran~\ref{app:ui-mockup}.

\subsection{Layar Tinjauan Hasil AI}

Layar tinjauan AI menampilkan instansi prediksi teratas, peringkat tiga prediksi teratas, \textit{bar confidence} per kelas, serta instansi tujuan terdekat yang dikembalikan backend. Pengguna dapat mengoreksi kategori secara manual sebelum laporan disimpan. Layar ini menjadi titik intervensi kunci antara model AI dan keputusan manusia, sehingga sinyal ketidakpastian ditonjolkan. Tampilan layar dapat dilihat pada Lampiran~\ref{app:ui-mockup}.

\subsection{Layar Rekomendasi Instansi}

Layar rekomendasi instansi menampilkan hasil routing dari backend berupa nama instansi tujuan, kategori instansi, dan estimasi jarak dari koordinat laporan. Informasi ini berasal dari prediksi 9 instansi dan koordinat latitude-longitude, bukan dari agregasi ke kluster baru. Tampilan layar dapat dilihat pada Lampiran~\ref{app:ui-mockup}.

\subsection{Layar Riwayat Laporan}

Layar riwayat menampilkan daftar laporan milik pengguna dalam bentuk kartu yang memuat \textit{thumbnail} foto, kategori instansi, tanggal pengiriman, dan status terkini. Pengguna dapat memilih satu kartu untuk melihat detail lengkap, melacak audit pada tabel \texttt{report\_history}, dan menerima pembaruan secara \textit{realtime} melalui Supabase \textit{Realtime}. Pada peran petugas, layar detail juga menyediakan aksi perubahan status; status selesai hanya dapat dikirim setelah foto bukti penyelesaian diunggah. Tampilan layar dapat dilihat pada Lampiran~\ref{app:ui-mockup}.

\subsection{Layar Peta Laporan}

Layar peta menampilkan pin lokasi seluruh laporan pengguna pada peta interaktif. Setiap pin dapat diketuk untuk memunculkan \textit{popup info} ringkas berisi kategori, status, dan tanggal. Filter kategori memungkinkan pengguna menyaring tampilan berdasarkan sembilan instansi target. Tampilan layar dapat dilihat pada Lampiran~\ref{app:ui-mockup}.

\subsection{Layar Profil dan Pengaturan}

Layar profil menampilkan informasi akun pengguna (foto, nama, surel, dan peran) yang diambil dari tabel \texttt{profiles} di Supabase. Pengaturan yang tersedia meliputi pengelolaan sesi, preferensi bahasa, dan opsi keluar yang membersihkan sesi Supabase dan \textit{state} Riverpod. Tampilan layar dapat dilihat pada Lampiran~\ref{app:ui-mockup}.

% =====================================================================
\section{Skenario Pengujian dan Metrik Evaluasi}
\label{sec:skenario-pengujian}
% =====================================================================

Bagian ini merangkum skenario pengujian yang dijalankan untuk memvalidasi pipeline klasifikasi multimodal beserta metrik evaluasi yang dipakai untuk mengukur kinerja sistem.

\subsection{Pembagian Data}

Dataset Jakarta CRM dipartisi secara stratifikasi berdasarkan label kategori, dengan proporsi 70 persen \textit{train}, 15 persen validasi, dan 15 persen uji. Partisi akhir berisi 43.241 sampel pelatihan, 9.266 sampel validasi, dan 9.266 sampel uji. Ketidakseimbangan kelas dikompensasi melalui pembobotan kelas (\textit{class weights}) pada Persamaan~\ref{eq:classweight}, bukan melalui \textit{resampling}. Pembagian disimpan pada berkas \texttt{train.csv}, \texttt{val.csv}, dan \texttt{test.csv} pada folder \texttt{artifacts/splits/} agar setiap eksperimen dapat menggunakan partisi yang identik.

\subsection{Metrik Evaluasi}

Metrik evaluasi yang dilaporkan terdiri atas dua kelompok, yaitu metrik utama dan metrik pendukung.

\textbf{Metrik utama:} \textit{macro F1 score} dipilih sebagai metrik utama karena memberikan bobot yang sama pada setiap kelas, sesuai untuk dataset dengan ketidakseimbangan kelas yang tidak ekstrem namun tetap signifikan. \textit{Balanced accuracy} dilaporkan sebagai metrik pendamping yang juga memberikan bobot yang sama pada setiap kelas. \textit{Confusion matrix} digunakan untuk menampilkan distribusi prediksi benar dan salah pada setiap pasangan kelas.

\textbf{Metrik pendukung:} \textit{precision} dan \textit{recall} per kelas dilaporkan untuk mengidentifikasi kelas mana yang mengalami kesulitan tertentu. Untuk konfigurasi yang menggunakan PGS, \textit{mean uncertainty} dihitung sebagai indikator kualitas kalibrasi probabilitas. Latensi inferensi end-to-end pada server FastAPI dihitung sebagai jumlah dari latensi pra-proses, latensi \textit{image encoder}, latensi \textit{text encoder}, latensi \textit{classifier head}, dan latensi pasca-proses.

\subsection{Skenario Ablasi}

Ablasi yang dilaporkan pada Bab~\ref{ch:hasil} terdiri atas tiga skenario, yaitu (a) ablasi modalitas dengan tiga konfigurasi (citra saja, teks saja, dan fusi awal), (b) ablasi kalibrasi dengan dua konfigurasi (CatBoost tanpa PGS dan CatBoost dengan PGS), serta (c) ablasi pasangan \textit{encoder} dengan enam belas kombinasi pada matrix grid lengkap (kombinasi empat \textit{image encoder}, termasuk DINOv3 \textit{large}, dengan empat \textit{text encoder}). Hasil setiap skenario disajikan dengan metrik utama dan metrik pendukung yang sama.

\subsection{Validasi Paritas ONNX}

Setelah pelatihan selesai, model \textit{native} CatBoost diekspor ke format ONNX. Validasi paritas dilakukan dengan menjalankan kedua model pada batch sampel uji yang sama, lalu menghitung selisih probabilitas absolut maksimum ($|\Delta p|_{\max}$) di seluruh komponen vektor keluaran. Kriteria yang dipakai adalah $|\Delta p|_{\max} < 10^{-5}$ pada sampel set uji, yang memastikan prediksi top-1 identik antara jalur PyTorch dan jalur ONNX.

\subsection{Pengujian Aplikasi Flutter}

Aplikasi Flutter diuji secara end-to-end pada perangkat Android fisik dengan tiga skenario, yaitu (a) alur pelaporan oleh Warga Pelapor mulai dari pengambilan foto, penulisan narasi, klasifikasi AI, tinjauan, hingga penyimpanan ke Supabase, (b) alur verifikasi oleh Operator Instansi pada antrean laporan yang sudah dikategorisasi, dan (c) alur moderasi lintas instansi oleh Moderator. Setiap skenario dievaluasi dari sisi waktu respons, konsumsi data jaringan, dan ketepatan klasifikasi yang ditampilkan pada layar tinjauan AI.

\subsection{Catatan tentang Pemilihan Algoritma}

Penelitian ini memilih CatBoost sebagai satu-satunya algoritma \textit{classifier head} pada seluruh skenario eksperimen, dengan justifikasi teknis dan operasional yang telah dipaparkan pada Bagian~\ref{sec:justifikasi-catboost}. Tidak ada perbandingan dengan algoritma klasifikasi lain pada Bab~\ref{ch:hasil}; ablasi yang dilaporkan adalah ablasi modalitas, ablasi kalibrasi, dan ablasi pasangan \textit{encoder}.

\bigskip

Bab ini telah memaparkan metodologi penelitian secara terstruktur, mulai dari kerangka berpikir, analisis kebutuhan dengan tiga peran pengguna, akuisisi dan persiapan dataset Jakarta CRM melalui \textit{web scraping}, lingkungan komputasi RunPod yang diakses melalui SSH, pra-proses citra dan teks, ekstraksi \textit{embedding} dari empat \textit{encoder} citra dan empat \textit{encoder} teks, fusi awal berbasis konkatenasi, justifikasi pemilihan CatBoost sebagai \textit{classifier head} dan kalibrasinya dengan PGS, rancangan eksperimen ablasi, metrik evaluasi, ekspor ONNX, perancangan server FastAPI dan Docker untuk inferensi server-side, perancangan aplikasi Flutter SmartCityApps beserta UML lengkap (\textit{use case}, \textit{activity}, \textit{class}, dan \textit{sequence diagram}), mekanisme perutean instansi terdekat berdasarkan latitude-longitude, bukti penyelesaian laporan, rancangan layar, hingga skenario pengujian dan metrik evaluasi yang akan dipakai pada Bab~\ref{ch:hasil}. Bab Hasil dan Pembahasan berikutnya akan melaporkan capaian metrik kuantitatif, analisis kesalahan, dan validasi end-to-end pada aplikasi.
